\scanpage{29}
with a marked weight \deleted{$\unsure{\nu}$}, \unsure{(}and of course, a given genus). From the point of view of graph structure associated with the numerical type $(g,\nu)$, it is the simplest of all! We shall study the most complex ones later on, which correspond to the case of a zero moduli dimension, \unsure{ie|\deleted{\unsure{or}}and} of a maximal moduli codimension, equal to $3(g-1)+\nu$.

It is tempting to introduce the compactification
\[
\widehat{\mathcal M}_G\overset{\mathrm{definition}}{=}\prod_{\alpha\in S}\widehat{\mathcal M}_{g_\alpha,\widehat I_\alpha}
\]
where the $\widehat{\mathcal M}_{g_\alpha,\widehat I_\alpha}$ are the Mumford–Deligne compactifications of the $\mathcal M_{g_\alpha,\widehat I_\alpha}\simeq\mathcal M_{g_\alpha,\widehat\nu_\alpha}$. \deleted{\unsure{They|It}}\inserted{For fixed $\alpha$, $\widehat{\mathcal M}_{g_\alpha,\widehat I_\alpha}$} represents\deleted{\unsure{nt}} the fibered groupoid
\[
(30)\qquad\begin{aligned}
\mathcal S\longmapsto&\ \text{relative Mumford–Deligne curves }\inserted{Y_\alpha}\text{ over }\mathcal S,\\
&\text{of type }g_\alpha,\widehat I_\alpha\quad\text{i.e. of genus}\\
&g_\alpha,\text{ equipped with a monomorphism}\\
&(\widehat I_\alpha)_{\mathcal S}\hookrightarrow Y_\alpha^{\text{smooth}}
\end{aligned}
\]
so that $\widehat{\mathcal M}_G$ represents the fibered groupoid
\mspage{28}
